\newpage
\subsection{Résultats principaux}

\textbf{\textit{Manipulation 1 : Ondes centimétriques -- Micro-ondes}}

Cette manipulation avait pour objectif d'étudier les propriétés des micro-ondes émises par un cornet de 9,4~GHz. Les tensions ont été lues au multimètre à la sortie du récepteur, pour différentes positions, angles et obstacles.

Les résultats quantitatifs principaux sont les suivants. Pour la réflexion, l'angle de réflexion correspond à l'angle d'incidence pour les quatre angles testés, de 30\textdegree{} à 60\textdegree. L'écart maximal, obtenu à 50\textdegree, vaut
\[ \varepsilon_{\max} = 4{,}0~\% \]
ce qui est inférieur à l'incertitude relative de 4,1\,\%, la loi de la réflexion est validée aux incertitudes expérimentales.

Pour les ondes stationnaires, la fréquence de l'émetteur déduite de la longueur d'onde vaut
\[ f = (8{,}8 \pm 0{,}5)~\text{GHz} \]
soit un écart de 6,1\,\% avec la valeur nominale de 9,4~GHz, à la limite de l'incertitude relative de 5,9\,\%. Une seule paire minimum-maximum ayant été relevée, plusieurs nœuds successifs et une régression linéaire rendraient ce résultat plus fiable.

Pour la réfraction, l'indice de la paraffine, calculé à partir de la déviation minimale, vaut
\[ n = 1{,}219 \pm 0{,}006 \]
soit un écart d'environ 19\,\% avec comme valeur attendue 1,5.

Le \refit{Tableau}{\ref}{tab:resume-qualitatif} résume les résultats des autres expériences, qui sont qualitatives.

\begin{table}[H]
    \centering
    \begin{tabular}{l p{4.6cm} p{5.6cm}}
        \toprule
        Expérience & Mesure principale & Observation \\
        \midrule
        1 -- Propagation & $U_{\max} = 332$~mV dans l'axe & Signal maximal lorsque les cornets sont alignés \\
        2 -- Pénétration & $U = 1{,}06$~V derrière le bois (8~cm) & Le bois laisse passer une partie du signal \\
        3 -- Blindage & $U = 0$~V derrière l'aluminium & L'aluminium bloque totalement le signal \\
        7 -- Huygens & $U$ de 54,45~mV à 112~mV avec la fente & La fente renvoie du signal dans la zone d'ombre \\
        8 -- Diffraction & $U_{\min} = 45$~mV derrière la plaque & Du signal subsiste dans la zone d'ombre \\
        9 -- Interférence & $U_{\max} = 0{,}424$~V au centre & Alternance de maxima et de minima \\
        10 -- Polarisation & 2,605~V (vertical), 0,002~V (horizontal) & Les micro-ondes émises sont polarisées \\
        \bottomrule
    \end{tabular}
    \caption{Tableau récapitulatif, résultats principaux, manipulation 1}
    \label{tab:resume-qualitatif}
\end{table}

\medskip

\textbf{\textit{Manipulation 2 : Trombone de Koenig -- Vitesse du son dans les gaz}}

Cette manipulation avait pour objectif de mesurer la vitesse du son dans l'air et dans le CO$_2$ avec le trombone de Koenig. Le déplacement $\Delta x$ du tube mobile a été relevé entre plusieurs maxima, pour différentes fréquences, et la vitesse a été obtenue par régression linéaire passant par l'origine.

Les résultats quantitatifs principaux sont les suivants. Pour l'air, à $T = 297{,}1$~K, la vitesse du son mesurée sur sept fréquences, de $1\,500$ à $3\,000$~Hz, vaut
\[ v_{\text{exp}} = (335{,}3 \pm 2{,}0)~\text{m/s} \]
soit un écart de 2,9\,\% ($10{,}0$~m/s) avec la valeur théorique de $(345{,}3 \pm 0{,}9)$~m/s, supérieur à l'incertitude combinée de $2{,}2$~m/s.

Pour le CO$_2$, à $T = 296{,}3$~K, la vitesse du son mesurée sur onze fréquences, de $1\,000$ à $2\,000$~Hz, vaut
\[ v_{\text{exp}} = (272{,}2 \pm 1{,}5)~\text{m/s} \]
soit un écart de 1,5\,\% ($3{,}9$~m/s) avec la valeur théorique de $(268{,}3 \pm 0{,}7)$~m/s, supérieur à l'incertitude combinée de $1{,}7$~m/s. Les causes supposées sont une présence résiduelle d'air dans le tube et le petit nombre de maxima ($p = 2$ ou $3$).

Le \refit{Tableau}{\ref}{tab:resume-trombone} résume les résultats pour les deux gaz.

\begin{table}[H]
    \centering
    \begin{tabular}{l p{4.6cm} p{5.6cm}}
        \toprule
        Gaz & Mesure principale & Observation \\
        \midrule
        Air & $v_{\text{exp}} = (335{,}3 \pm 2{,}0)$~m/s, écart de 2,9\,\% & $\Delta x$ proportionnel à $p/(2f)$, vitesse indépendante de la fréquence \\
        CO$_2$ & $v_{\text{exp}} = (272{,}2 \pm 1{,}5)$~m/s, écart de 1,5\,\% & $\Delta x$ proportionnel à $p/(2f)$, vitesse indépendante de la fréquence \\
        Hélium & -- & Non mesuré, le gaz n'étant pas disponible \\
        \bottomrule
    \end{tabular}
    \caption{Tableau récapitulatif, résultats principaux, manipulation 2}
    \label{tab:resume-trombone}
\end{table}

La loi théorique de la vitesse du son est validée aux incertitudes expérimentales, avec des écarts inférieurs à 3\,\%.
